\section*{\texorpdfstring{\S\CurrentExerciseGroup}{第{\CurrentExerciseGroup}节}}
\phantomsection
\addcontentsline{toc}{section}{第{\CurrentExerciseGroup}节习题}

\begin{enumerate}
\def\labelenumi{\arabic{enumi})}
\item
  设 \((\lambda_{\alpha})\) 是 T 上的复点测度组成的一个族。为使族 \((\lambda_{\alpha})\) 在赋有淡拓扑的向量空间 \(\mathcal{M}(\mathrm{T};\mathbf{C})\) 中可和，必须且只须族 \((|\lambda_{\alpha}|)\) 可和。
\item
  设 \((\lambda_{n})\) 是 T 上的复测度序列。为使序列 \((\lambda_{n})\) 在赋有淡拓扑的 \(\mathcal{M}(\mathrm{T};\mathbf{C})\) 中可和，必须且只须对每个函数 \(f\in \mathcal{K}(\mathrm{T})\) ，通项为 \(\lambda_{n}(f)\) 的级数绝对收敛（利用 Banach-Steinhaus 定理；参见 TVS, III, §4, No.~2, 定理 1 的推论 2）。
\item
  设 \(\mathbf{T}\) 是紧区间 \([0,1]\) （\(\mathbf{R}\) 的紧区间）；对每个整数 \(n > 0\) ，设 \(\lambda_{n}\) 是测度
\end{enumerate}

\[
f \mapsto \frac {1}{n} \int_ {0} ^ {1} f (x) \sin 2 n \pi x d x
\]

（T 上的）。证明：在赋有淡拓扑（或拟强拓扑（Ch. III, §1, 习题 8））的空间 \(\mathcal{M}(\mathrm{T})\) 中，测度序列 \((\lambda_n)\) 可和，但序列 \((|\lambda_n|)\) 不可和。（注意若 \(\alpha_n = n \cdot \lambda_n(f)\) 则 \(\sum_{n} \alpha_n^2 < +\infty\) ，这要利用 Bessel 不等式，再由 Cauchy-Schwarz 不等式由此推出 \(\sum_{n} |\lambda_n(f)| < +\infty\) 。）

\begin{enumerate}
\def\labelenumi{\arabic{enumi})}
\setcounter{enumi}{3}
\tightlist
\item
  设 \(\mu\) 是 T 上的一个正测度，\((\nu_{i})_{i\in I}\) 是 T 上的正测度组成的可和族，使得 \(\mu \leqslant \sum_{i}\nu_{i}\) 。存在 T 上的测度组成的可和族 \((\mu_{i})_{i\in I}\) ，使得 \(\mu = \sum_{i}\mu_{i}\) 且 \(\mu_{i} \leqslant \nu_{i}\) （对所有 \(i \in I\) ）。（利用命题 4，化归到 T 紧的情形，再化归到 \(I = N\) 的情形。然后用归纳法构造诸 \(\mu_{i}\) 。）
\end{enumerate}

\setcounter{exercisegroup}{3}
\edef\CurrentExerciseGroup{\arabic{exercisegroup}}
